\section{Proposed IB-CellFree-GNN Framework}

The proposed IB-CellFree-GNN framework fundamentally shifts the paradigm of cell-free MIMO receiver design from a traditional decoupled architecture---where signal quantization and multi-user detection are optimized independently---to a joint Quantization-Aware Training (QAT) approach. Theoretically grounded in the Variational Information Bottleneck (VIB) principle, the framework aims to compress the local observations at the APs such that only the most task-relevant features for downstream multi-user detection are preserved, thereby minimizing backhaul overhead. By jointly optimizing the distributed quantization policy at the APs and the centralized graph-based interference cancellation at the CPU, the system dynamically adapts to varying channel conditions and strict backhaul capacity constraints.

To enable end-to-end training of the discrete quantization process, we introduce a differentiable Learned Step Size Quantization (LSQ) module implemented via custom autograd functions in PyTorch. The quantizer employs the Straight-Through Estimator (STE) to bypass the non-differentiability of the rounding operation during backpropagation, allowing the quantization step size to be optimized via gradient descent. Building upon this, we design a progressive residual quantization structure that categorizes the AP observations into distinct precision levels: 0-bit (complete dropping of deeply faded signals), 2-bit (extraction of the most significant bits for basic reliability), and 4-bit (fine-grained refinement on the residuals for high-precision recovery). Mathematically, the forward pass of the quantizer scales the input by a learnable step size, clamps it to the target bit-width range, rounds it to the nearest integer, and scales it back, while the STE ensures that the gradients flow smoothly to the preceding layers.

At the core of the distributed processing is the Resource-Inference Policy Network, deployed locally at each AP. This bit-width controller is implemented as a lightweight Multi-Layer Perceptron (MLP) with a 3-32-3 architecture that takes the concatenated local soft symbols and the local signal-to-noise ratio (SNR) as inputs. Crucially, to bridge the discrete bit-selection operation with the continuous optimization required by deep learning, the policy network utilizes the Gumbel-Softmax trick. During the training phase, this technique generates a differentiable, continuous approximation of the categorical distribution over the available bit-widths, allowing gradients to backpropagate through the sampling process. As the training progresses, an annealing temperature parameter $\tau$ is gradually decayed, smoothly transitioning the soft assignments into hard categorical choices. During the deployment phase, the network executes a strict hard-decision ($\arg\max$) inference, ensuring that the APs transmit exactly the inferred number of bits without any floating-point overhead.

At the CPU, a bit-aware heterogeneous Graph Neural Network (GNN) is designed to aggregate the variable-length feature vectors received from the distributed APs. The GNN first performs a feature lifting process, mapping the concatenated quantized symbols, local SNR, and the reported bit-depth (expected\_bits) into a high-dimensional hidden space. An AP-to-CPU attention mechanism is then applied to dynamically adjust the spatial aggregation weights based on the reported bit-depth, inherently suppressing the impact of heavily quantized or dropped signals. Subsequently, a Transformer-based Multi-User Interference Cancellation (MU-IC) module processes the fused features to resolve spatial interference among users. The entire framework is trained end-to-end using a joint loss function that elegantly combines a task-oriented Mean Squared Error (MSE) loss for signal detection with a bandwidth penalty term formulated as $\lambda \times \text{expected\_bits}$. This joint formulation, operating in tandem with the annealing temperature parameter $\tau$, ensures that the system learns to achieve near-optimal performance while strictly adhering to the available backhaul capacity.

A rigorous theoretical analysis of the proposed IB-CellFree-GNN reveals its exceptional scalability and efficiency. The computational complexity of the traditional local LMMSE preprocessing at each AP scales as $\mathcal{O}(N^3 + N^2 K)$ due to the matrix inversion. In contrast, the deployed Resource-Inference Policy Network processes the signals user-by-user using a fixed-size MLP, resulting in a computational complexity of $\mathcal{O}(K)$ per AP. Consequently, the AI-driven decision module exhibits an $\mathcal{O}(1)$ computational complexity scaling with respect to the number of AP antennas $N$. At the CPU, the bit-aware GNN and Transformer-based MU-IC operate on the aggregated spatial features with a complexity of $\mathcal{O}(K^2 d + K d^2)$, where $d$ is the hidden dimension. This centralized processing complexity is entirely independent of $L$ and $N$ after the initial spatial pooling. Furthermore, the Gumbel-Softmax annealing strategy provides empirical convergence guarantees by ensuring that the gradient variance remains bounded as the continuous relaxation asymptotically approaches the discrete categorical distribution.